\section{Estimate D$^\dagger$: opening, linearization, and the
two-frequency reduction}\label{sec:reduction}

Throughout this section $q\sim Q=x^{u'}$ and $p\sim P=x^{u}$ are band
primes, $u,u'\in(1/3,1/2)$, and the $n$-side configuration is $n=pa$ with
$n\equiv1\pmod q$, so that $a$ ranges over $a\le x/p$,
$a\equiv\bar p\pmod q$. We write $R:=x/(pq)$ for the progression length
and $\Llog=\log x$.

\subsection{Opening the progression: primitive characters mod $q^2$}

\begin{lemma}[$\beta$; AP-opening]\label{lem:beta}
\statusPM\quad
Let $q$ be an odd prime, $\lambda\not\equiv0\pmod q$, and $w$ any weight.
Then
\[
\sum_{\substack{n\le x\\ n\equiv1\,(q)}}w(n)\chi_\lambda(n)
=\frac{1}{\varphi(q)}\sum_{\chi_1\ \mathrm{mod}\ q}\
\sum_{\substack{n\le x\\(n,q)=1}}w(n)\,\big(\tilde\chi_1\chi_\lambda\big)(n),
\]
where $\tilde\chi_1$ is the lift of $\chi_1$ to modulus $q^2$; and
\emph{every} product $\tilde\chi_1\chi_\lambda$ is primitive of conductor
exactly $q^2$.
\end{lemma}

\begin{proof}
The identity is orthogonality of the characters mod $q$ applied to
$\mathbf 1[n\equiv1\ (q)]$ on $(n,q)=1$. For primitivity: a character mod
$q^2$ is imprimitive iff it is trivial on the kernel
$1+q\Z/q^2$ of reduction $(\Z/q^2)^*\to(\Z/q)^*$. The lift
$\tilde\chi_1$ is trivial there by construction, while
$\chi_\lambda(1+qu)=\e{-\lambda u/q}$ by \eqref{eq:fq-hom} is a nontrivial
character of $1+q\Z/q^2\cong\Z/q$ when $\lambda\ne0$. Hence the product is
nontrivial on the kernel and has conductor $q^2$.
\end{proof}

\begin{remark}[the assembly accounting; a negative result]\label{rem:negative}
Lemma~$\beta$ composed with averaged bounds for primitive characters to
square moduli (Lemma~$\alpha$, Appendix~\ref{app:alpha}) does \emph{not}
prove Estimate~D$^\dagger$: the opening converts a log-saving target on a
length-$x/q$ object into a power-of-$q$-saving target on length-$x$
objects, while every family-average tool returns square-root-type savings
--- short by a factor $\approx q^{1-\kappa}$ (maximal-character route) or
$x^{2u'-1/2}$ (full-family bilinear route). The deeper reason: the
load-bearing hypothesis of the averaged machinery is $q$-independence of
the weight, and D$^\dagger$'s weight is $q$-dependent \emph{through the
progression itself}, $m\mapsto w(1+qm)$ --- essentially, not technically.
This dead end is recorded (with the two bracketing Cauchy--Schwarz routes)
so it is not re-entered; the lemmas $\alpha,\beta$ are kept as tools.
\end{remark}

\subsection{Linearization on the progression}

\begin{lemma}[linearization]\label{lem:linearization}
\statusPM\quad
Let $(a_0,q)=1$ and $a=a_0+qt$. Then for every $\lambda$,
\[
\chi_\lambda(a_0+qt)\;=\;\chi_\lambda(a_0)\,
\e{-\frac{\lambda\,\bar a_0\,t}{q}},
\qquad \bar a_0a_0\equiv1\ (\mathrm{mod}\ q).
\]
In particular, on $a\equiv\bar p\pmod q$ one has $\bar a_0\equiv p$, and
the deep character becomes the \emph{classical} additive character
$\e{-\lambda pt/q}$ in the progression coordinate $t$, with frequency
$\lambda p/q$ moving with $p$.
\end{lemma}

\begin{proof}
$a_0+qt\equiv a_0(1+qt\bar a_0)\pmod{q^2}$ (expand; the error is
$q^2t\bar a_0\cdot(\ldots)$). Apply \eqref{eq:fq-hom}:
$\fq(a_0+qt)=\fq(a_0)+\fq(1+qt\bar a_0)=\fq(a_0)-t\bar a_0\bmod q$.
\end{proof}

\begin{remark}
Machine-checked to $10^{-14}$ against the true Fermat-quotient phase
(\texttt{wp21\_test.py}). Consequence: for $v\equiv1$ the inner sum is
geometric, $|T_p(\lambda)|\le\min(R,\nm{\lambda p/q}^{-1})$ --- the
unweighted layer of D$^\dagger$ is main-term bookkeeping. All difficulty
lives in the decorated weight.
\end{remark}

\subsection{Unconditional coverage of the deep harmonics}

The binding digit condition at $q$ contributes Fourier weights
$c_\lambda$ (of the interval $[q/2,q)$ in $\Z/q$) with
$|c_\lambda|\ll\langle\lambda\rangle^{-1}$, where
$\langle\lambda\rangle:=\min(|\lambda|,q-|\lambda|)$. Set
\[
\Lambda_0\;:=\;x^{\,u+2u'-1+\varepsilon}
\qquad(\text{note }1<u+2u'<\tfrac32\text{ in the band, and }\Lambda_0<q).
\]

\begin{proposition}[large $\lambda$ are covered]\label{prop:largelambda}
\statusP\quad
Let $v$ be any $1$-bounded weight supported on
$\{a\le x/p:\ a\equiv\bar p\ (\mathrm{mod}\ q)\}$ (no structure assumed).
Then for every band prime $p$,
\[
\sum_{\Lambda_0\le\langle\lambda\rangle<q/2}|c_\lambda|\,
\Big|\sum_{a}v(a)\chi_\lambda(a)\Big|
\;\ll\;\frac{x}{pq}\,x^{-\varepsilon/3},
\]
and summing over $p\sim P$ bounds the deep-harmonic part of
$E_q$ by $(x/q)\,x^{-\varepsilon/3}$: a power saving, far below the
$L^{-c}$ requirement.
\end{proposition}

\begin{proof}
Write $T_p(\lambda)=\sum_av(a)\chi_\lambda(a)$ and let
$N_a=x/p$ be the range length, so the support size is
$\le R+1=x/(pq)+1$ and $\|v\|_2^2\le R+1$. Decompose the $\lambda$-range
into dyadic blocks $\langle\lambda\rangle\sim\Lambda$,
$\Lambda_0\le\Lambda<q/2$. By Cauchy--Schwarz and
Corollary~\ref{cor:ap} (progression-restricted deep large sieve, with
$c=\bar p$),
\[
\sum_{\langle\lambda\rangle\sim\Lambda}|c_\lambda||T_p(\lambda)|
\;\le\;\Big(\sum_{\langle\lambda\rangle\sim\Lambda}|c_\lambda|^2\Big)^{1/2}
\Big(\sum_{\lambda\bmod q}|T_p(\lambda)|^2\Big)^{1/2}
\;\ll\;\Lambda^{-1/2}\Big(\frac{N_a}{q}+q\Big)^{1/2}(R+1)^{1/2}.
\]
Since $N_a/q=R$ and $q/R=x^{u+2u'-1}$, the right side is
\[
\ll\;\Lambda^{-1/2}R\,\Big(\frac1R+\frac{q}{R^2}\Big)^{1/2}\sqrt{2}
\;\ll\;R\,\Big(\frac{x^{u+2u'-1}}{\Lambda}\Big)^{1/2}
\;\le\;R\,x^{-\varepsilon/2}
\]
for $\Lambda\ge\Lambda_0$ (using $q/R^2\ge1/R$ in the band). Summing over
the $O(\log x)$ dyadic blocks loses a log, absorbed by
$x^{-\varepsilon/2}\log x\ll x^{-\varepsilon/3}$.
\end{proof}

\begin{remark}
Only $0<\langle\lambda\rangle<\Lambda_0$ --- polynomially many harmonics,
but crucially \emph{including the bounded ones} --- survive. By a
standard completion/Vaaler step the remaining task is equivalent to
Estimate~D$^\dagger$ for each fixed $|\lambda|\le\Llog^{C}$; the empirics
of Appendix~\ref{app:numerics} measure exactly these objects and find
square-root cancellation with room to spare.
\end{remark}

\subsection{The two-frequency reduction of the digit layer}

Insert the $n$-side decoration. The binding $n$-side digit condition at
$p$ is $\fr{n/p^2}\ge\tfrac12$, i.e.\ $a\bmod p\in[p/2,p)$, detected by
$f:\Z/p\to\R$, the \emph{centered} indicator of $[p/2,p)$, with Fourier
expansion $f(y)=\sum_{0<\mu<p}\hat f_\mu\e{\mu y/p}$,
$\hat f_0=0$, $|\hat f_\mu|\le\tfrac1{2\langle\mu\rangle}$,
$\langle\mu\rangle=\min(\mu,p-\mu)$. The cofactor-anatomy indicator is
$\sigma$ ($1$-bounded; \S\ref{sec:anatomy} handles it; this subsection
takes $\sigma\equiv1$).

\begin{proposition}[two-frequency reduction]\label{prop:two-freq}
\statusPM\quad
Fix the band primes $p,q$ and $0<|\lambda|\le\Llog^C$, and let
$a_0\equiv\bar p\ (\mathrm{mod}\ q)$ parametrize the progression
$a=a_0+qt$, $0\le t\le R$. Then the digit-decorated inner object of
Estimate~D$^\dagger$ equals
\[
V_p(\lambda)\;:=\;\sum_{t\le R}\sigma(a_0+qt)\,
f\big((a_0+qt)\bmod p\big)\,\e{-\frac{\lambda pt}{q}}
\;=\;\sum_{0<\mu<p}\hat f_\mu\,\e{\frac{\mu a_0}{p}}
\sum_{t\le R}\sigma(a_0+qt)\,\e{t\,\theta_\mu},
\]
with the \emph{two-frequency phases}
\[
\theta_\mu\;:=\;\frac{\mu q}{p}-\frac{\lambda p}{q}
\;=\;\frac{\mu q^2-\lambda p^2}{pq}\pmod 1 .
\]
For $\sigma\equiv1$ each inner sum is geometric:
$\big|\sum_{t\le R}\e{t\theta_\mu}\big|\le\min\big(R,\tfrac1{2\nm{\theta_\mu}}\big)$.
\end{proposition}

\begin{proof}
Substitute the Fourier expansion of $f$ and
$\e{\mu(a_0+qt)/p}=\e{\mu a_0/p}\e{(\mu q/p)t}$, then combine with the
linearized deep phase $\e{-(\lambda p/q)t}$ of
Lemma~\ref{lem:linearization}. The geometric bound is standard.
\end{proof}

\begin{definition}[the dichotomy]\label{def:dichotomy}
The digit layer of D$^\dagger$ is a family of linear exponential sums at
the phases $\theta_\mu$. Call $p$ \emph{minor} for $(q,\lambda)$ when no
small $\mu$ makes $\nm{\theta_\mu}$ abnormally small and $q/p$ is not
abnormally well approximable below scale $R$
(Definition~\ref{def:minor-conditions} makes this precise); otherwise
$p$ is \emph{major}. Since
$\theta_\mu\approx0\iff p/q\approx\sqrt{\mu/\lambda}$, the major $p$
localize on explicit Diophantine curves; they are counted, not bounded,
in \S\ref{sec:major}. The flagged empirical drift at small
$\lambda$ and top of band (Appendix~\ref{app:numerics}) is exactly the
$\mu/\lambda=1$ curve ($p\approx q$) announcing itself: structure to be
counted, not noise.
\end{definition}
